\HMTNarrativeHeading{Qué acopla la constante}

La constante holográfica de acoplamiento aritmético-geométrico se introduce
donde se produce su registro. Su representación integral se denota por
\(K\); la lectura periódica, por \(\kappa_{\rm ag}\). El origen de
fase, la orientación y el orden de las componentes forman parte de la
definición. La palabra «constante» nombra el objeto determinado por esta
construcción y por sus lectores, que se expondrán mediante operaciones
explícitas.

La extracción reúne los canales orientados y un dato de suma. La inversión
de Hadamard recompone sus componentes en el orden natural. Cada
coeficiente de esa inversión procede del operador declarado y de su
imagen integral. De este modo, el registro queda determinado antes de
utilizarse en la clausura de la constante de estructura fina. La lectura
posicional convierte después las doce componentes en una expresión
racional periódica. Fijadas la base y la longitud, esta conversión admite
una inversa exacta que restituye todos los bloques.

La misma ordenación tiene una función geométrica. Sus lectores distinguen
una cuaterna marcada y una configuración de incidencia; la firma asociada
a alfa selecciona una hexada dentro del diseño correspondiente. La
relación con las regiones numéricas utiliza las aplicaciones de
alineamiento y elevación que se demostrarán en la parte excepcional.
La lectura periódica y la configuración incidencial reciben así una
fuente común, con operaciones de distinto tipo.

El desarrollo operatorio posterior añade una tercera función. Los
desplazamientos cíclicos del registro forman una base de su portador. Las
respuestas sobre esa base determinan un operador, y la forma de Gram
proporciona una cota de estabilidad para recuperarlo. Cuando existe la
simetría cíclica especificada, una respuesta vectorial completa permite
reconstruir las demás. El archivo isométrico conserva esa capacidad de
identificación.

Éste es el contenido del acoplamiento que vertebra la exposición:
codificación aritmética reversible, incidencia marcada e identificación
operatoria. Primero se construye el registro; después se prueba lo que
cada lector obtiene de él. La teoría de recuperación se desarrollará
cuando estén disponibles sus espacios y transportes. Esta posición
explica conjuntamente la centralidad de la constante y la jerarquía de
las demostraciones que utilizan sus diferentes realizaciones.

